% Vietnamese translation for OI Public Library
% Source-PDF: IOI中国国家候选队论文/2009/徐持衡/浅谈几类背包题.pdf
\documentclass[11pt]{article}
\usepackage{oipl-vn}

\newcommand{\Oh}{\mathcal{O}}

\title{Bàn sơ lược về vài dạng bài ba lô}
\author{Xu Chiheng\\Trường Trung học Ôn Châu, Chiết Giang\\Người hướng dẫn: Shu Chunping}
\date{Tháng 12 năm 2008}

\begin{document}
\maketitle

\origtitle{Bàn sơ lược về vài dạng bài ba lô}
\sourcepath{IOI China candidate team papers / 2009 / Xu Chiheng / knapsack variants}

\begin{abstract}
Bài toán ba lô là một bài toán kinh điển và là một phần rất cơ bản trong quy
hoạch động. Tuy nhiên, lấy bài toán ba lô \(0\)-\(1\) làm gốc, người ta có thể
biến đổi ra rất nhiều dạng đề khác nhau; lời giải của chúng tất nhiên cũng
không giống nhau, nên rất đáng tiếp tục khảo sát.

Bài viết này thảo luận bốn bài toán biến thể của ba lô và trình bày một số
cách làm của tác giả, với mong muốn gợi mở thêm hướng suy nghĩ.
\end{abstract}

\noindent\textbf{Từ khóa:} quy hoạch động, ba lô, tối ưu hóa.

\section{Mở đầu}

Bài toán ba lô là một bài toán tối ưu kinh điển trong vận trù học. Nó là một
bài toán NP-đầy đủ, nhưng lại có bối cảnh ứng dụng thực tế rất rộng, xuất phát
từ một tình huống quen thuộc trong đời sống:
một chiếc ba lô và rất nhiều vật phẩm; ta cần bỏ một số vật phẩm vào ba lô để
đạt một mục tiêu nhất định.

Trong tin học, sau khi lượng hóa mọi dữ liệu, ta thu được bài toán sau.

\term{Bài toán ba lô \(0\)-\(1\).}
Cho \(n\) vật phẩm và một ba lô. Vật phẩm \(i\) có giá trị \(W_i\), thể tích
\(V_i\), ba lô có dung lượng \(C\). Có thể tùy ý chọn các vật phẩm bỏ vào ba
lô; hãy tìm tổng giá trị lớn nhất của các vật phẩm được bỏ vào.

Khi chọn vật phẩm cho vào ba lô, đối với mỗi vật phẩm \(i\), hoặc bỏ nó vào,
hoặc không bỏ. Không được bỏ vật phẩm \(i\) nhiều lần, cũng không được chỉ bỏ
một phần của vật phẩm \(i\) (tức là không được chia vật phẩm). Vì vậy bài toán
này được gọi là bài toán ba lô \(0\)-\(1\).

Các phương pháp thường dùng để giải ba lô \(0\)-\(1\) gồm quay lui, tham lam,
thuật toán di truyền, tìm kiếm tabu, mô phỏng luyện kim, v.v.

Ở bậc phổ thông, cái gọi là bài toán ba lô \(0\)-\(1\) kinh điển thường đảm
bảo mọi dữ liệu sau khi lượng hóa đều là số nguyên dương, tức là một trường
hợp đặc biệt của quy hoạch nguyên. Trong bài viết này, nếu không nói khác, ta
đều dùng giả thiết đó. Lời giải quy hoạch động kinh điển có độ phức tạp
\(\Oh(nC)\) như sau.

Trạng thái: xét \(i\) vật phẩm đầu tiên, chọn một số vật phẩm sao cho tổng thể
tích không vượt quá \(j\), giá trị lớn nhất đạt được là \(F_i[j]\). Quyết định
hiện tại là vật phẩm thứ \(i\) có được chọn hay không. Ta có chuyển tiếp:
\[
  F_i[j]=F_{i-1}[j]\qquad (V_i>j\ge 0),
\]
\[
  F_i[j]=\max\{F_{i-1}[j],\,F_{i-1}[j-V_i]+W_i\}
  \qquad (C\ge j\ge V_i).
\]

Do \(F_i\) chỉ liên quan tới \(F_{i-1}\), có thể dùng mảng cuốn để tiết kiệm
bộ nhớ. Pseudocode kinh điển là:

\begin{verbatim}
FOR i := 1 TO n
    FOR j := C DOWNTO Vi
        Max(F[j], F[j - Vi] + Wi) -> F[j]
    END FOR
END FOR
\end{verbatim}

\section{Các biến đổi cơ bản của ba lô}

\subsection{Ba lô hoàn toàn}

\term{Bài toán ba lô hoàn toàn.}
Cho \(n\) loại vật phẩm và một ba lô. Loại vật phẩm thứ \(i\) có giá trị
\(W_i\), thể tích \(V_i\), ba lô có dung lượng \(C\), và số lượng mỗi loại vật
phẩm là vô hạn. Có thể tùy ý chọn các vật phẩm bỏ vào ba lô; hãy tìm tổng giá
trị lớn nhất của các vật phẩm được bỏ vào.

Bài này hoàn toàn có thể chuyển thành ba lô \(0\)-\(1\): chia loại vật phẩm
thứ \(i\) thành \(\lfloor C/V_i\rfloor\) vật phẩm riêng biệt, rồi áp dụng quy
hoạch động kinh điển của ba lô \(0\)-\(1\). Tuy nhiên độ phức tạp thời gian
của cách này quá cao, không phải một cách cài đặt thực dụng.

Dễ thấy bài toán này là một cực khác so với ba lô \(0\)-\(1\): mỗi loại vật
phẩm đều có thể lấy không giới hạn. Chỉ cần đổi phương trình chuyển tiếp là có
thể xây dựng thuật toán mới.

Trạng thái: xét \(i\) loại vật phẩm đầu tiên, chọn một số vật phẩm sao cho tổng
thể tích không vượt quá \(j\), giá trị lớn nhất đạt được là \(F_i[j]\). Quyết
định hiện tại là có chọn loại vật phẩm thứ \(i\) hay không, và nếu chọn thì có
thể chọn nhiều cái. Chuyển tiếp là:
\[
  F_i[j]=F_{i-1}[j]\qquad (V_i>j\ge 0),
\]
\[
  F_i[j]=\max\{F_{i-1}[j],\,F_i[j-V_i]+W_i\}
  \qquad (C\ge j\ge V_i).
\]
Chú ý hạng thứ hai là \(F_i\), không phải \(F_{i-1}\). Điểm khác ba lô
\(0\)-\(1\) chỉ nằm ở đây, vì ta được phép thêm một vật phẩm nữa trên cơ sở đã
chọn trước đó.

Như vậy bài toán có lời giải \(\Oh(nC)\), cùng bậc với ba lô \(0\)-\(1\), và
cũng có thể dùng mảng cuốn:

\begin{verbatim}
FOR i := 1 TO n
    FOR j := Vi TO C
        Max(F[j], F[j - Vi] + Wi) -> F[j]
    END FOR
END FOR
\end{verbatim}

\subsection{Ba lô nhiều lần}

\term{Bài toán ba lô nhiều lần.}
Cho \(n\) loại vật phẩm và một ba lô. Loại vật phẩm thứ \(i\) có giá trị
\(W_i\), thể tích \(V_i\), số lượng \(K_i\), ba lô có dung lượng \(C\). Có thể
tùy ý chọn các vật phẩm bỏ vào ba lô; hãy tìm tổng giá trị lớn nhất của các
vật phẩm được bỏ vào.

Giống ba lô hoàn toàn, bài này cũng có thể trực tiếp dùng quy hoạch động kinh
điển của ba lô \(0\)-\(1\), nhưng hiệu quả quá thấp, nên cần tìm thuật toán
hiệu quả hơn.

Trước hết, với loại vật phẩm thứ \(i\), ta không thể xác định ngay lấy bao
nhiêu cái là tối ưu, vì không có gì chứng minh rằng lấy một cái hoặc lấy hết
sẽ nhất định tốt hơn. Nói cách khác, số lượng cần trong nghiệm tối ưu có thể
là bất cứ số nào từ \(0\) đến \(K_i\).

Trong đời sống, nếu cần có khả năng đưa ra bất kỳ số tiền nguyên nào từ \(1\)
đến \(K_i\), ta thường không mang \(K_i\) tờ một đơn vị, vì quá bất tiện; thay
vào đó ta mang vài tờ mệnh giá \(1\) và một số tờ có mệnh giá khác nhau. Ý
tưởng này có thể dùng trực tiếp trong bài toán.

Không cần chia một loại vật phẩm thành \(K_i\) vật phẩm đơn lẻ; chỉ cần chia
nó thành các nhóm sao cho có thể ghép được mọi số lượng từ \(1\) đến \(K_i\).
Có thể chứng minh rằng cách tách nhị phân làm số nhóm nhỏ nhất: tách \(K_i\)
thành
\[
  1,2,4,\ldots,2^t,\;K_i-2^{t+1}+1
  \qquad (2^{t+2}>K_i\ge 2^{t+1}).
\]
Như vậy đã đủ đáp ứng mọi số lượng có thể xuất hiện trong nghiệm tối ưu. Sau
đó vẫn dùng thuật toán kinh điển của ba lô \(0\)-\(1\). Ta thu được thuật toán
có độ phức tạp
\[
  \Oh\!\left(C\sum_i \lfloor \log_2 K_i\rfloor\right).
\]

\begin{verbatim}
FOR i := 1 TO n
    1 -> m
    WHILE Ki > 0
        IF m > Ki THEN Ki -> m
        Ki - m -> Ki
        FOR j := C DOWNTO Vi * m
            Max(F[j], F[j - Vi * m] + Wi * m) -> F[j]
        END FOR
        m * 2 -> m
    WEND
END FOR
\end{verbatim}

Thuật toán này còn có một tối ưu: nếu \(K_i>C/V_i\), thì phần vượt quá là vô
nghĩa và có thể bỏ đi. Độ phức tạp khi đó có thể viết là
\(\Oh(nC\log_2 C)\). Vì trong các bài ở bậc phổ thông, giá trị \(C\) thường
có giới hạn vừa phải, thuật toán này trong thực tế đã đủ đáp ứng nhu cầu
thông thường.

Tuy vậy, phần sau có một cách giải hiệu quả hơn.

\subsection{Tối ưu bằng hàng đợi đơn điệu}

\term{Bài toán tổng đoạn liên tiếp lớn nhất có giới hạn độ dài.}
Cho dãy \(X_i\) dài \(n\). Hãy tìm tổng lớn nhất của một đoạn liên tiếp có độ
dài không vượt quá \(l_{\max}\).

Trước hết xét cách đơn giản nhất: dùng hai vòng lặp \(\Oh(nl_{\max})\) để tìm
giá trị lớn nhất.

\begin{verbatim}
FOR i := 1 TO n
    0 -> s
    FOR j := i DOWNTO Max(i - lmax + 1, 1)
        s + Xj -> s
        IF s > ans THEN s -> ans
    END FOR
END FOR
\end{verbatim}

Gọi \(S_i\) là tổng từ \(X_1\) đến \(X_i\). Nếu đã cố định một đầu mút, việc
còn lại là lấy một cực trị trên một đoạn liên tiếp của mảng \(S\). Bài toán
cực trị đoạn có thể giải bằng cây đoạn. Nhưng bài này còn có tính chất khác:
độ dài đoạn là cố định theo cửa sổ trượt, và mỗi đoạn chỉ cần truy vấn một lần.
Vì vậy có thể dùng cấu trúc đơn giản hơn: hàng đợi đơn điệu.

Xét hàng đợi đơn điệu để duy trì giá trị lớn nhất. Sau khi bảo đảm chỉ số trong
dãy tăng dần, ta còn yêu cầu giá trị phần tử giảm dần:
\[
\begin{array}{c|ccccc}
\text{vị trí trong hàng đợi} & L & L+1 & \cdots & R-1 & R\\
\hline
\text{chỉ số} & A_L & A_{L+1} & \cdots & A_{R-1} & A_R\\
\text{giá trị} & B_L & B_{L+1} & \cdots & B_{R-1} & B_R
\end{array}
\]
Thỏa mãn \(A_{i+1}>A_i>A_{i-1}\) và \(B_{i-1}>B_i>B_{i+1}\) với
\(R>i>L\).

Ngoài thao tác lấy phần tử đầu hàng đợi, hàng đợi đơn điệu cần thêm thao tác
để duy trì tính chất đặc biệt của nó. Nếu đưa vào một phần tử mới \((a,b)\),
thì \(a\) chắc chắn lớn hơn \(A_R\), nhưng \(b\) có thể lớn hơn hoặc bằng
\(B_R\). Vì \(b\ge B_R\), trong khi phần tử \(R\) sẽ rời hàng đợi trước phần
tử mới, nên phần tử \(R\) đã mất giá trị sử dụng: từ nay về sau phần tử mới
chắc chắn không kém phần tử \(R\). Do đó ta có thể xóa phần tử \(R\).

Lặp lại bước trên cho đến khi phía trước không còn phần tử nào hoặc thỏa mãn
\(B_R>b\), rồi mới cho phần tử mới vào hàng đợi. Rõ ràng phần tử đầu hàng đợi
hiện tại chính là giá trị lớn nhất của đoạn.

\begin{verbatim}
PROCEDURE INSERT a, b
    WHILE R >= L AND b > B[R] DO R - 1 -> R
    R + 1 -> R
    a -> A[R]
    b -> B[R]
END
\end{verbatim}

Hàng đợi đơn điệu như vậy không bảo đảm mỗi thao tác riêng lẻ đều là \(\Oh(1)\),
nhưng vì mỗi phần tử chỉ vào hàng đợi một lần và ra hàng đợi một lần, tổng
hiệu quả là \(\Oh(n)\).

Dĩ nhiên, bài toán vừa nêu chỉ là chức năng cơ bản của hàng đợi đơn điệu. Điều
ta cần là dùng nó để tối ưu ba lô, nên ta xét lại bài toán ba lô nhiều lần.

Với loại vật phẩm thứ \(i\), biết thể tích \(v\), giá trị \(w\), số lượng
\(k\). Có thể chia toàn bộ mảng quy hoạch động thành \(v\) phần theo phần dư
của thể tích \(j\) modulo \(v\). Ví dụ với \(v=3\):
\[
\begin{array}{c|cccccccccc}
j & 0&1&2&3&4&5&6&7&8&\cdots\\
\hline
j\bmod v & 0&1&2&0&1&2&0&1&2&\cdots
\end{array}
\]

Ta xử lý riêng từng phần. Giả sử phần dư là \(d\):
\[
\begin{array}{c|ccccccc}
\text{chỉ số }j & 0&1&2&3&4&5&\cdots\\
\hline
\text{thể tích tương ứng}
& d&d+v&d+2v&d+3v&d+4v&d+5v&\cdots
\end{array}
\]

Sau khi chia nhóm, trạng thái có chỉ số \(j\) có thể chuyển từ bất kỳ chỉ số
nào trong đoạn \(j-k\) đến \(j-1\), vì hai thể tích kề nhau đúng bằng \(v\).
Điều này đã khá giống bài toán cực trị đoạn. Nhưng cần chú ý rằng do thể tích
khác nhau, giá trị ở thể tích lớn hiển nhiên có thể lớn hơn hoặc bằng thể tích
nhỏ; so sánh trực tiếp là vô nghĩa. Vì vậy cần hiệu chỉnh giá trị về cùng một
mốc thể tích. Chẳng hạn cùng quy về \(d\), tức là đưa
\[
  F[jv+d]-j w
\]
vào hàng đợi thay cho giá trị ban đầu.

Với vật phẩm thứ \(i\), pseudocode chuyển tiếp là:

\begin{verbatim}
FOR d := 0 TO v - 1                 // duyet phan du, xu ly rieng
    xoa rong hang doi
    FOR j := 0 TO (C - d) div v     // j la chi so, the tich la j*v+d
        INSERT j, F[j * v + d] - j * w
        IF A[L] < j - k THEN L + 1 -> L
        B[L] + j * w -> F[j * v + d]
    END FOR
END FOR
\end{verbatim}

Vì hàng đợi đơn điệu có tổng độ phức tạp tuyến tính, ba lô nhiều lần sau tối
ưu bằng hàng đợi đơn điệu có độ phức tạp \(\Oh(nC)\).

\section{Một vài dạng ba lô khác}

\subsection{Ba lô phụ thuộc dạng cây: chọn môn}

\term{Bài toán ba lô phụ thuộc dạng cây.}
Cho \(n\) vật phẩm và một ba lô. Vật phẩm thứ \(i\) có giá trị \(W_i\), thể
tích \(V_i\), nhưng phụ thuộc vào vật phẩm thứ \(X_i\): chỉ sau khi chọn
\(X_i\) mới được chọn \(i\); nếu không có phụ thuộc thì \(X_i=0\). Quan hệ phụ
thuộc tạo thành một rừng. Ba lô có dung lượng \(C\). Có thể tùy ý chọn vật
phẩm bỏ vào ba lô; hãy tìm tổng giá trị lớn nhất của các vật phẩm được bỏ vào.

Bài này cần dùng ba lô trên nền tảng quy hoạch động trên cây.

\subsubsection*{Vật phẩm tổng quát}

Xét một loại ``vật phẩm'' không có chi phí (thể tích) và giá trị cố định; thay
vào đó, giá trị của nó thay đổi theo chi phí (thể tích) mà ta phân bổ cho nó.

Vật phẩm tổng quát có thể biểu diễn bằng một mảng một chiều mô tả quan hệ giữa
thể tích và giá trị: \(G_j\) là giá trị tương ứng khi thể tích là \(j\), với
\(C\ge j\ge 0\). Rõ ràng mảng quy hoạch động ba lô \(F_i\) trước đây chính là
một vật phẩm tổng quát, vì \(F_i[j]\) biểu diễn giá trị lớn nhất khi thể tích
là \(j\). Tương tự, nhiều vật phẩm cũng có thể gộp thành một vật phẩm tổng
quát.

\subsubsection*{Tổng của vật phẩm tổng quát}

Phép gộp hai vật phẩm tổng quát thành một vật phẩm tổng quát chính là việc liệt
kê cách phân bổ thể tích cho hai vật phẩm, thỏa mãn:
\[
  G[j]=\max_{0\le k\le j}\{G_1[j-k]+G_2[k]\}
  \qquad (C\ge j\ge 0).
\]

Thời gian gộp hai vật phẩm tổng quát là \(\Oh(C^2)\).

Với bài toán ba lô có quan hệ phụ thuộc dạng cây, có thể xem mỗi cây con là
một vật phẩm tổng quát. Khi đó vật phẩm tổng quát của một cây con bằng tổng
của vật phẩm tổng quát ứng với gốc cây con và các vật phẩm tổng quát ứng với
mọi cây con nối với gốc.

Toàn bộ quá trình quy hoạch động đi từ lá lên gốc. Với mỗi cây con, thao tác
là:

\begin{verbatim}
PROCEDURE DEAL i, v, w       // nut i, the tich v, gia tri w
    FOR j := v TO C          // khoi tao vat pham tong quat nay
        w -> Fi[j]
    END FOR
    FOR s := 1 TO n
        IF s la con cua i THEN
            tong cua Fi va Fs -> Fi
        END IF
    END FOR
END
\end{verbatim}

Mỗi lần chỉ gộp được hai vật phẩm tổng quát. Muốn gộp \(n\) vật phẩm tổng quát
thành một vật phẩm cần \(n-1\) lần gộp, nên thuật toán có hiệu quả
\(\Oh(nC^2+n^2)\). Phần \(\Oh(n^2)\) tất nhiên có thể dùng danh sách kề để
giảm thành \(\Oh(n)\), do đó độ phức tạp cuối cùng là \(\Oh(nC^2)\).

Nhìn lại ba lô \(0\)-\(1\) kinh điển: trong thuật toán đó, tổng của một vật
phẩm tổng quát và một vật phẩm thường chỉ cần \(\Oh(C)\). Từ đó suy ra:

\subsubsection*{Tổng của vật phẩm tổng quát và một vật phẩm thường}

Gộp một vật phẩm tổng quát với một vật phẩm thường thành một vật phẩm tổng
quát có thể làm bằng cách tương tự quy hoạch động kinh điển của ba lô
\(0\)-\(1\):
\[
  G[j]=G_1[j]\qquad (v>j\ge 0),
\]
\[
  G[j]=\max\{G_1[j],\,G_1[j-v]+w\}
  \qquad (C\ge j\ge v).
\]

Việc gộp như vậy chỉ tốn \(\Oh(C)\), cũng là gộp thêm một vật phẩm nhưng nhanh
hơn rất nhiều so với tổng của hai vật phẩm tổng quát. Vậy có cách nào dùng
kiểu gộp \(\Oh(C)\) này thay cho việc tính tổng của hai vật phẩm tổng quát để
xử lý bài toán hay không?

\subsubsection*{Hợp của vật phẩm tổng quát}

Vì giữa hai vật phẩm tổng quát có phần giao, nên không thể đồng thời chọn cả
hai. Khi đó ta cần lấy \term{hợp} của vật phẩm tổng quát: với cùng một thể
tích, chọn cái có giá trị lớn hơn, độ phức tạp \(\Oh(C)\):
\[
  G[j]=\max\{G_1[j],\,G_2[j]\}\qquad (C\ge j\ge 0).
\]

Xét lại việc xử lý cây con gốc \(i\). Giả sử hiện cần xử lý một con \(s\) của
\(i\). Nếu trong \(F_i\) hiện tại ta bắt buộc đưa vật phẩm \(s\) vào làm trạng
thái ban đầu của cây con gốc \(s\), thì sau khi xử lý xong cây con gốc \(s\),
\(F_s\) là một vật phẩm tổng quát có giao với \(F_i\) (thực tế là \(F_s\) chứa
\(F_i\)). Đồng thời, \(F_s\) phải thỏa mãn đã chọn vật phẩm \(s\), nên
\(F_s[j]\) với \(V_s>j\ge 0\) đã vô nghĩa, còn \(F_s[j]\) với \(C\ge j\ge V_s\)
chắc chắn chứa vật phẩm \(s\). Để tiện cài đặt, sau khi xử lý, chương trình
quy định chỉ \(F_s[j]\) với \(C-V_s\ge j\ge 0\) là hợp lệ.

Tiếp theo chỉ cần lấy hợp của \(F_s\) và \(F_i\) rồi gán cho \(F_i\), là hoàn
tất việc xử lý một con. Như vậy tổng độ phức tạp chỉ là \(\Oh(nC)\).

\begin{verbatim}
PROCEDURE DEAL i, C
    FOR s := 1 TO n
        IF s la con cua i THEN
            Fi -> Fs
            DEAL s, C - Vs
            FOR k := Vs TO C
                Max(Fi[k], Fs[k - Vs] + Ws) -> Fi[k]
            END FOR
        END IF
    END FOR
END
\end{verbatim}

Dùng thuật toán này có thể tối ưu các bài kiểu ``Chọn môn'' xuống \(\Oh(nC)\).
Nó cũng có thể dùng làm một lời giải quy hoạch động trên cây \(\Oh(nC)\) cho
bài ``Kế hoạch ngân sách của Jinming'' của NOIP 2006. Sau khi bài viết hoàn
thành, tác giả được TKY nhắc rằng thuật toán cùng độ phức tạp đã được đề cập
trong bài năm 2007 ``Bàn về tổ chức dữ liệu hợp lý'' của He Sen; vì vậy cách
làm ở đây chỉ nên xem là một cài đặt \(\Oh(nC)\) đơn giản hơn.

\subsection{PKU 3093}

\term{PKU 3093.}
Cho \(n\) vật phẩm và một ba lô. Vật phẩm thứ \(i\) có thể tích \(V_i\), ba lô
có dung lượng \(C\). Yêu cầu chọn một số vật phẩm bỏ vào ba lô sao cho mọi vật
phẩm còn lại đều không thể bỏ thêm vào được; hãy tính số phương án.

Tạm thời chưa xét điều kiện ``mọi vật phẩm còn lại đều không thể bỏ thêm vào
được''. Khi đó bài toán chỉ là đếm mọi phương án khả thi của ba lô \(0\)-\(1\).

Dùng \(F_i[j]\) biểu diễn số phương án chọn một số vật phẩm trong \(i\) vật
phẩm đầu tiên sao cho tổng thể tích bằng \(j\). Khởi tạo \(F_0[0]=1\). Chuyển
tiếp:
\[
  F_i[j]=F_{i-1}[j]\qquad (V_i>j),
\]
\[
  F_i[j]=F_{i-1}[j]+F_{i-1}[j-V_i]\qquad (j\ge V_i).
\]
Rõ ràng thuật toán này có hiệu quả \(\Oh(nC)\), và nó tính số phương án của
mọi cách đặt vật phẩm vào ba lô.

Bây giờ xét điều kiện ``mọi vật phẩm còn lại đều không thể bỏ thêm vào''. Nếu
vật phẩm nhỏ nhất trong các vật phẩm còn lại có thể tích \(v\), thì số phương
án là
\[
  \sum_{C\ge j>C-v} F_n[j],
\]
với điều kiện ta đã xác định trước vật phẩm nhỏ nhất còn lại là vật phẩm nào.

Sau khi sắp xếp theo thể tích, bước tiếp theo là liệt kê \(i\) làm một vật
phẩm nhỏ nhất trong phần còn lại. Với mọi vật phẩm \(s<i\), bắt buộc phải bỏ
vào ba lô; với \(i\), không được bỏ vào; với mọi vật phẩm \(s>i\), làm thống
kê phương án khả thi của ba lô \(0\)-\(1\), rồi cộng
\(\sum F_n[j]\) với \(C\ge j>C-V_i\) vào \(\mathit{ans}\). Vì mỗi lần cần
thống kê trên \(n-i\) vật phẩm và tổng cộng thống kê \(n\) lần, độ phức tạp là
\(\Oh(n^2C)\).

\begin{verbatim}
0 -> sum                         // sum la tong the tich vat pham 1..i-1
FOR i := 1 TO N
    xoa mang F
    1 -> F[sum]                  // khoi tao
    FOR s := i + 1 TO N          // thong ke phuong an kha thi
        FOR j := C DOWNTO Vs + sum
            F[j] + F[j - Vs] -> F[j]
        END FOR
    END FOR
    FOR k := C DOWNTO C - Vi + 1 // cong vao tong so phuong an
        IF k >= sum THEN ans + F[k] -> ans
    END FOR
    sum + Vi -> sum
END FOR
\end{verbatim}

Có thể thấy cùng một vật phẩm được xét đưa vào ba lô nhiều lần, gây lãng phí
thời gian. Quan sát thấy vật phẩm thứ \(i\) được xét tổng cộng \(i-1\) lần.
Mỗi vòng lặp lại bớt đi một vật phẩm. Nếu đảo ngược toàn bộ quá trình, liệu
mỗi vật phẩm có thể chỉ xét một lần không?

Do trạng thái khởi tạo khác nhau, trước hết cần thống nhất trạng thái khởi
tạo. Có thể để mỗi lần \(F[0]=1\), tổng dung lượng là \(C-\mathit{sum}\).
Nhưng chỉ thống nhất khởi tạo vẫn chưa đủ, vì dung lượng ba lô mỗi lần vẫn
khác nhau. Khi làm thống kê ba lô, dung lượng không được là một biến thay đổi,
mà cũng phải thống nhất; vì vậy mỗi lần xét một vật phẩm đều dùng dung lượng
lớn nhất \(C\) để cập nhật ba lô.

Sau một thao tác, cộng
\[
  \sum F[j]\qquad
  (C-\mathit{sum}\ge j>C-\mathit{sum}-V_i,\; j\ge 0)
\]
vào \(\mathit{ans}\). Khi đó mỗi vật phẩm chỉ được xét một lần, dung lượng ba
lô thống nhất là \(C\), và độ phức tạp trở thành \(\Oh(nC)\).

\begin{verbatim}
0 -> sum
1 -> F[0]
FOR i := 1 TO n
    sum + Vi -> sum
END FOR
FOR i := n DOWNTO 1
    sum - Vi -> sum
    FOR j := C - sum DOWNTO Max(C - sum - Vi + 1, 0)
        ans + F[j] -> ans
    END FOR
    FOR j := C DOWNTO Vi
        F[j] + F[j - Vi] -> F[j]
    END FOR
END FOR
\end{verbatim}

\section{Tổng kết}

Nhìn lại bốn bài toán ba lô trong bài viết: với ba lô hoàn toàn, chỉ cần đổi
phương trình chuyển tiếp; với ba lô nhiều lần, dùng tối ưu hàng đợi đơn điệu
để đạt \(\Oh(nC)\); trong ba lô phụ thuộc dạng cây, ta đưa vào khái niệm mới
và cuối cùng chuyển hóa thuật toán; còn ở bài PKU 3093, ta hợp nhất các thao
tác tương tự để tối ưu.

Mặc dù các phương pháp được dùng không giống nhau, mỗi phương pháp đều không
chỉ giới hạn trong bài toán ba lô, mà hoàn toàn có thể vận dụng linh hoạt ở
những bài toán khác.

Mọi bài được nhắc tới trong bài đều được giải bằng thuật toán thời gian
\(\Oh(nC)\). Điều này không có nghĩa mọi bài ba lô đều có thể tối ưu đến mức
đó, nhưng có thể khẳng định rằng không thể có độ phức tạp nhanh hơn tuyến tính
theo kích thước bảng \(nC\) nếu vẫn phải xử lý toàn bộ bề mặt trạng thái này.

Hiện nay vẫn còn rất nhiều bài toán dạng ba lô chưa có lời giải thật tốt, chờ
mọi người tiếp tục tìm tòi và nghiên cứu.

\appendix
\section{Tài liệu tham khảo}

\begin{itemize}
  \item \textit{Chín bài giảng về ba lô}, DDengi, ZJU.
\end{itemize}

\section{Các đề bài gốc được nhắc trong bài}

\subsection{Chọn môn}

Nguồn: CTSC 1997.

Ở đại học áp dụng chế độ tín chỉ. Mỗi môn học có một số tín chỉ nhất định; chỉ
cần sinh viên chọn môn đó và qua kiểm tra thì có thể nhận số tín chỉ tương
ứng. Tổng tín chỉ cuối cùng của sinh viên là tổng tín chỉ của các môn đã chọn.

Mỗi sinh viên đều phải chọn một số lượng môn quy định. Một số môn có thể chọn
trực tiếp, còn một số môn cần kiến thức nền tảng nhất định: chỉ sau khi đã
chọn một số môn khác mới được chọn. Ví dụ, ``Cấu trúc dữ liệu'' phải được chọn
sau ``Lập trình ngôn ngữ bậc cao''; khi đó ta gọi môn sau là môn tiên quyết
của môn trước. Mỗi môn có nhiều nhất một môn tiên quyết trực tiếp. Hai môn học
cũng có thể có cùng một môn tiên quyết. Để tiện mô tả, mỗi môn có một mã số
liên tiếp \(1,2,3,\ldots\). Ví dụ:

\[
\begin{array}{c|c|c}
\text{mã môn} & \text{môn tiên quyết} & \text{tín chỉ}\\
\hline
1 & \text{không có} & 1\\
2 & 1 & 1\\
3 & 2 & 3\\
4 & \text{không có} & 3\\
5 & 2 & 4
\end{array}
\]

Trong ví dụ trên, \(1\) là môn tiên quyết của \(2\); nếu muốn chọn \(2\) thì
nhất định \(1\) đã được chọn. Tương tự, nếu muốn chọn \(3\), thì \(1\) và \(2\)
đều phải đã được chọn.

Sinh viên không thể học hết mọi môn mà trường đại học mở, nên phải chọn trước
những môn mình muốn học khi nhập học. Tổng số môn mỗi sinh viên được chọn là
cho trước. Hãy tìm một phương án chọn môn sao cho nhận được nhiều tín chỉ nhất
và phải thỏa mãn nguyên tắc môn tiên quyết đi trước. Giả sử giữa các môn không
có xung đột thời gian.

\paragraph{Input.}
Dòng đầu của tệp vào gồm hai số nguyên dương \(M,N\) cách nhau bởi một dấu
cách; \(M\) là tổng số môn chờ chọn \((1\le M\le 1000)\), \(N\) là tổng số môn
sinh viên có thể chọn \((1\le N\le M)\).

Tiếp theo là \(M\) dòng, mỗi dòng mô tả một môn, mã môn lần lượt là
\(1,2,\ldots,M\). Mỗi dòng có hai số: số đầu là mã môn tiên quyết của môn này
\((0\) nếu không có\()\), số thứ hai là tín chỉ. Tín chỉ là số nguyên dương
không vượt quá \(10\).

\paragraph{Output.}
Dòng đầu chỉ có một số, là tổng tín chỉ của các môn thực sự được chọn. Tiếp
theo là \(N\) dòng, mỗi dòng một số, biểu diễn mã môn được chọn.

\paragraph{Ví dụ.}
\begin{verbatim}
INPUT.TXT
7 4
2 2
0 1
0 4
2 1
7 1
7 6
2 2

OUTPUT.TXT
13
2
6
7
3
\end{verbatim}

\paragraph{Mã nguồn của bài ``Chọn môn''.}
\begin{verbatim}
var
    n,C,i:longint;
    x,w:array[1..400]of longint;
    f:array[0..400,0..400]of longint;

procedure dfs(k,C:longint);
var
    i,j:longint;
begin
    if C<=0 then exit;
    for i:=1 to n do if x[i]=k then begin
        for j:=0 to C-1 do f[i,j]:=f[k,j]+w[i]; {bat buoc chon i}
        dfs(i,C-1);
        for j:=1 to C do
            if f[i,j-1]>f[k,j] then
                f[k,j]:=f[i,j-1]; {lay hop}
    end;
end;

begin
    readln(n,C);
    for i:=1 to n do begin
        readln(x[i],w[i]); {cha x[i] va tin chi w[i]}
    end;
    dfs(0,C); {0 la cha cua moi nut khong co cha}
    writeln(f[0,C]);
end.
\end{verbatim}

\subsection{Kế hoạch ngân sách của Jinming}

Nguồn: NOIP 2006, bài 2.

\paragraph{Mô tả bài toán.}
Jinming rất vui vì căn nhà mới của gia đình sắp được nhận chìa khóa; trong nhà
mới có một căn phòng rộng rãi dành riêng cho cậu. Điều làm cậu vui hơn là hôm
qua mẹ nói: ``Phòng của con cần mua những đồ gì, bố trí thế nào thì con tự
quyết, miễn là không vượt quá \(N\) đồng.'' Sáng nay Jinming bắt đầu lập ngân
sách. Cậu chia các món muốn mua thành hai loại: vật chính và phụ kiện. Phụ
kiện thuộc về một vật chính nào đó, ví dụ:

\[
\begin{array}{c|l}
\text{vật chính} & \text{phụ kiện}\\
\hline
\text{máy tính} & \text{máy in, máy quét}\\
\text{tủ sách} & \text{sách}\\
\text{bàn học} & \text{đèn bàn, văn phòng phẩm}\\
\text{ghế làm việc} & \text{không có}
\end{array}
\]

Nếu muốn mua một vật được phân loại là phụ kiện, trước hết phải mua vật chính
mà phụ kiện đó thuộc về. Mỗi vật chính có \(0\), \(1\), hoặc \(2\) phụ kiện.
Phụ kiện không có phụ kiện của riêng nó. Jinming muốn mua rất nhiều thứ, chắc
chắn vượt quá giới hạn \(N\) đồng của mẹ. Vì vậy cậu quy định độ quan trọng
của từng món theo \(5\) mức, dùng số nguyên \(1\) đến \(5\), trong đó mức \(5\)
là quan trọng nhất. Cậu cũng tra trên Internet giá của từng món (đều là bội
của \(10\) đồng). Cậu muốn trong điều kiện không vượt quá \(N\) đồng, làm cho
tổng tích giữa giá tiền và độ quan trọng của mỗi món được mua là lớn nhất.

Giả sử món thứ \(j\) có giá \(v[j]\), độ quan trọng \(w[j]\), và đã chọn \(k\)
món với chỉ số \(j_1,j_2,\ldots,j_k\). Đại lượng cần tối đa là
\[
  v[j_1]w[j_1]+v[j_2]w[j_2]+\cdots+v[j_k]w[j_k].
\]

Hãy giúp Jinming thiết kế một danh sách mua sắm thỏa mãn yêu cầu.

\paragraph{Input.}
Dòng đầu của tệp \texttt{budget.in} gồm hai số nguyên dương cách nhau bởi dấu
cách:
\[
  N\quad m.
\]
Trong đó \(N<32000\) là tổng tiền, \(m<60\) là số món muốn mua. Từ dòng thứ
hai đến dòng \(m+1\), dòng thứ \(j\) cho dữ liệu cơ bản của món có số hiệu
\(j-1\), gồm ba số nguyên không âm:
\[
  v\quad p\quad q.
\]
Trong đó \(v<10000\) là giá, \(p\in\{1,\ldots,5\}\) là độ quan trọng, và \(q\)
cho biết vật này là vật chính hay phụ kiện: nếu \(q=0\), đây là vật chính; nếu
\(q>0\), đây là phụ kiện, và \(q\) là số hiệu của vật chính mà nó thuộc về.

\paragraph{Output.}
Tệp \texttt{budget.out} chỉ có một số nguyên dương, là giá trị lớn nhất không
vượt quá tổng tiền, nhỏ hơn \(200000\).

\paragraph{Ví dụ.}
\begin{verbatim}
Input
1000 5
800 2 0
400 5 1
300 5 1
400 3 0
500 2 0

Output
2200
\end{verbatim}

\paragraph{Chương trình trong bài gốc.}
\begin{verbatim}
var
    n,C,i:longint;
    x,w,v:array[1..60]of longint;
    f:array[0..60,0..3200]of longint;

procedure dfs(k,C:longint);
var i,j:longint;
begin
    if C<=0 then exit;
    for i:=1 to n do if x[i]=k then begin
        for j:=0 to C-v[i] do f[i,j]:=f[k,j]+w[i]; {bat buoc chon i}
        dfs(i,C-v[i]);
        for j:=v[i] to C do
            if f[i,j-v[i]]>f[k,j] then f[k,j]:=f[i,j-v[i]]; {lay hop}
    end;
end;

begin
    assign(input,'budget.in'); reset(input);
    assign(output,'budget.out'); rewrite(output);
    readln(C,n);
    C:=C div 10;
    for i:=1 to n do begin
        readln(v[i],w[i],x[i]); {chi phi, do quan trong, cha}
        w[i]:=w[i]*v[i];
        v[i]:=v[i] div 10;
    end;
    dfs(0,C); {0 la cha cua moi nut khong co cha}
    writeln(f[0,C]);
    close(output);
end.
\end{verbatim}

\subsection{Margaritas on the River Walk}

Nguồn: PKU 3093, Greater New York 2006.

\paragraph{Mô tả.}
Một hoạt động phổ biến ở San Antonio là thưởng thức margarita trong công viên
dọc con sông, khu vực được gọi là River Walk. Có thể mua margarita tại nhiều
cơ sở dọc River Walk, từ các khách sạn sang trọng đến quầy Joe's Taco and
Margarita. Giá margarita thay đổi tùy theo lượng và chất lượng nguyên liệu
cũng như không gian của cơ sở. Bạn dành ra một khoản tiền nhất định để thử
nhiều loại margarita khác nhau.

Cho giá của một ly margarita tại mỗi cơ sở và số tiền được phân bổ, hãy tính
có bao nhiêu tổ hợp cực đại khác nhau có thể mua, với điều kiện chọn nhiều
nhất một ly từ mỗi cơ sở. Một tổ hợp hợp lệ phải có tổng giá không vượt quá số
tiền được phân bổ, và số tiền chưa dùng phải nhỏ hơn giá của mọi cơ sở chưa
được chọn. Nếu không, ta vẫn còn có thể thêm cơ sở đó vào tổ hợp.

Ví dụ, giả sử có \(25\) đô-la và các giá nguyên đô-la là:
\[
\begin{array}{c|rrrrrr}
\text{Vendor} & A & B & C & D & H & J\\
\hline
\text{Price} & 8 & 9 & 8 & 7 & 16 & 5
\end{array}
\]

Các tổ hợp có thể là:
\[
\begin{gathered}
ABC(25),\ ABD(24),\ ABJ(22),\ ACD(23),\ ACJ(21),\\
ADJ(20),\ AH(24),\ BCD(24),\ BCJ(22),\ BDJ(21),\\
BH(25),\ CDJ(20),\ CH(24),\ DH(23),\ HJ(21).
\end{gathered}
\]
Vậy tổng số tổ hợp là \(15\).

\paragraph{Input.}
Dữ liệu bắt đầu bằng một dòng chứa số nguyên \(N\), là số bộ dữ liệu tiếp theo
\((1\le N\le 1000)\). Mỗi bộ dữ liệu bắt đầu bằng một dòng chứa hai số nguyên
\(V\) và \(D\), lần lượt là số cơ sở \((1\le V\le 30)\) và số đô-la được phép
chi \((1\le D\le 1000)\). Hai số cách nhau bởi một hoặc nhiều dấu cách. Phần
còn lại của mỗi bộ dữ liệu gồm một hoặc nhiều dòng, mỗi dòng chứa một hoặc
nhiều số nguyên là giá margarita của các cơ sở. Tổng cộng có \(V\) giá được
cho. Giá của một ly margarita luôn ít nhất là \(1\). Dữ liệu được chọn sao cho
kết quả nằm trong số nguyên không dấu \(32\) bit.

\paragraph{Output.}
Với mỗi bộ dữ liệu, in một dòng gồm số thứ tự bộ dữ liệu, một dấu cách, rồi số
tổ hợp của bộ dữ liệu đó.

\paragraph{Ví dụ.}
\begin{verbatim}
Sample Input
2
6 25
8 9 8 7 16 5
30 250
1 2 3 4 5 6 7 8 9 10 11
12 13 14 15 16 17 18 19 20
21 22 23 24 25 26 27 28 29 30

Sample Output
1 15
2 16509438
\end{verbatim}

\paragraph{Gợi ý.}
Một số phương pháp giải có thể có độ phức tạp lũy thừa theo số cơ sở. Với các
phương pháp này, giới hạn thời gian có thể bị vượt ở những bộ dữ liệu có số cơ
sở lớn, chẳng hạn ví dụ thứ hai ở trên.

\end{document}
